\documentclass[12px]{article}

\title{Lezione 16 Metodi Numerici di Approssimazione}
\date{2026-05-05}
\author{Federico De Sisti}

\input{../../setup.tex}

\begin{document}
	\maketitle
	\newpage
	\subsection{Metodi Multistep (Linear multistep method)}
	\begin{defi}[Metodo Multi Step]
	La soluzione al passo $n+1$ dipende non solo dalla soluzione al passo $n$ (ed eventualmente da se stessa nel caso di metodi impliciti) ma anche da $p+1$ passi precedenti, cioè dipende da $u_n,u_{n-1},u_{n-2}, \ldots,u_{n+1-p}$, \ \ \ $p\geq 0$.\\
	\end{defi}
	\textbf{Osservazione}\\
	Runge Kutta non è mutlistep, perchè calcola la soluzione in tempi intermedi tra $u_n$ e $u_{n+1}$.\\
	Diciamo $p+1$ perché quel $p$ torna comodo in futuro.\\
	Un metodo ad un passo ha $p=0$.\\
	L'approccio qui è furbo, si memorizza non solo le soluzioni ai tempi precedenti ma anche la dinamica calcolata ai tempi precedenti.\\
	Sono metodi della forma generale
	\[
		u_{n+1} = \sum^{p}_{j=0}a_ju_{n-j} + j \sum_{j={-1}}^{p}b_jf_{n-j} \ \ \ n\geq p
	.\] 
	dove $u_{n-j} \approx y(t_{n-j})$ e  $f_{n-j} = f(t_{n-j},u_{n-j})\approx f(t_{n-j}, y(t_{n-j}))$\\
	Se sostituisco quel  $j=-1$ noto che ottengo  $f_{n+1}$, quindi questa scrittura comprende anche i metodi impliciti.\\
	\textbf{Esempio} \\
	$p = 0 \Rightarrow  $ 1 passo $ \Rightarrow u_{n+1} = a_0u_n + h(b_{-1}f_{n+1}+b_0f_n)$ 
	\begin{itemize}
		\item $a_0 = 1 \ \ b_{-1} = 0 \ \ b_0 = 1 \Rightarrow  $ Eulero in avanti.
		\item $a_0 = 1 \ \  b_{-1} = 1 \ \ b_0 = 0 \Rightarrow  $ Eulero all'indietro.
		\item $a_0 = 1 \ \ b_{-1} = 1/2 \ \ b_0 = 1/2 \Rightarrow  $ Crank Nicholson.
	\end{itemize}
	L'obbiettivo è però quello di alzare l'ordine di convergenza.\\
	Teniamo in memoria\\
	$u_{n-p},\ldots,u_n$\\
	$f_{n-p}, \ldots, f_n, f_{n+1}$ se  $b_{-1}\neq 0$\\
	Una volta usati questi valori per calcolare  $u_{n+1}$, non mi serve più il valore  $u_{n-p}$ dato che utilizzo solamente i  $p+1$ passi precedenti, analogo per la dinamica $f_{n-p}$.\\
	Al contrario di Runge Kutta, che butta via i  $K_n$ calcolati ogni passo, qui "riciclo" quasi tutti i calcoli fatti.\\
\textbf{Osservazione}\\
Il metodo ha bisogno di $p+1$ dati precedenti, dove li prendo se ho solo un dato iniziale? Questi dati per iniziare si chiamano Valori di innesco.
\newpage
\subsubsection{Valori di Innesco}
Abbiamo bisogno di $u_0,u_1,\ldots,u_p$\\
$u_0$ è dato. Come calcolo gli altri valori? Potrei usare uno sviluppo di Taylor per i primi valori? Solitamente si calcolano con un metodo diverso, ad esempio Runge Kutta, oppure con un LMM con meno passi.\\
Potrei per esempio fare il primo passo con Eulero in avanti che per un passo ha ordine $2$, per poi proseguire con u metodo a 2 passi per esempio, questo posso farlo perché lo sviluppo di Taylor  genera un errore di ordine $h^4$
 \begin{itemize}
	 \item Se $b_{-1} = 0$ il metodo è esplicito
	 \item Se $b_{-1}= 0$ il metodo è implicito
\end{itemize}
\textbf{Esempi}\\
Se la prima sommatoria in LMM si riduce ad un solo elemento, cioè $\exists \bar J$ tale che  $a_{\bar J}=1$ e  tutti gli altri sono nulli
 \[
	 u_{n+1} = u_{n-\bar J} + h \sum^{p}_{j=-1}b_jf_{n-j} 
.\] 
Questa sarebbe una versione discreta di
\[
	y(t_{n+1}) = y(t_{n-\bar J}) + \int_{t_n-\bar J}^{t_n + 1}f(t,y(t))dt
.\] 
In cui $\int$ è approssimato da una formula di quadratura.\\
\textbf{Esempi}\\
$\bar J =1$ su  $[t_{n-1}, t_{n+1}] \Rightarrow \int_{t_{n-1}}^{t_{n+1}}f(t,y(t))dt\approx f(t_n,y(t_n))2h$
\[
	\Rightarrow u_{n+1} = u_{n-1} + 2hf_n
.\] 
\textbf{Esempi}\\
Simpson: $\bar J = 1$ su  $\displaystyle[t_{n-1},t_{n+1}] \Rightarrow  \int_{t_{n-1}}^{t_{n+1}}f(t,y(t))dt\approx \frac h 3(f_{n-1} + 4f_n + f_{n+1})$
\[
	\Rightarrow  u_{n+1} = u_{n-1} + \frac h 3 (f_{n-1} + 4f_n + f_{n+1})
.\] 
\textbf{Esempi}\\
$\bar J = p = 0 $ su  $[t_{n},t_{n+1}]$ + Rettangoli $ \Rightarrow $ Eulero in avanti
$\bar J = p = 0 $ su  $[t_{n},t_{n+1}]$ + Trapezi $ \Rightarrow $ Crank Nicholson\\
\hline\ \\
Se la seconda sommatoria in LMM si riduce al solo termine $n+1$
\[
\begin{split}
	u_{n+1} &= \sum^{p}_{j=0}a_ju_{n-j} + h b_{-1}f_{n+1}\\
	\frac{u_{n+1} - \sum^{p}_{j=0}a_j u_{n-j}}{hb_{-1}} &= f_{n+1}
\end{split}
.\] 
Con il quale stiamo approssimando $\dot y(t_{n+1}) = f(t_{n=1},y(t_{n+1},y_{n+1})$ per  $h \rightarrow 0$\\[10px]
In generale sembra quindi che la prima sommatoria sia legata alla derivata e la seconda sembra legata alla quadratura.\\
\textbf{Esempio}\\
$b_{-1} = 1$  \ \  $a_j = \begin{cases}
	1 \ \ j =0 \\
	0 \ \ \text{altrimenti}
\end{cases} \Rightarrow  $ Eulero all'indietro
\subsubsection{Consistenza di un LMM}
Come per i metodi ad un passo considero la soluzione esatta di ODE, la sostituisco nello schema e valuto il residuo.
\[
	h\tau_{n+1}(h) = y_{n+1} - \left[ \sum^{p}_{j=0}a_jy_{n-j} + h \sum^{p}_{j=-1}b_jf(t_{n-j},y_{n-j}) \right] \ \ \ n\geq p
.\] 
\[
	\tau(h) = \sup_{n} |\tau_{n+1}(h)|
.\] 
LMM è consistente se $\tau(h) \xrightarrow{h \rightarrow 0} 0$, di ordine $q$ se $\tau(h) = O(h^q)$\\
\begin{prop}
	LMM è consistente $ \Leftrightarrow$ $ \sum^{p}_{j=0}a_j = 1$ e $- \sum^{p}_{j=0}ja_j + \sum^{p}_{j=-1} b_j=-1$\\
	Inoltre se $y\in C^{q+1}((0,T)), \ q\geq 1$ allora LLM è consistente di ordine $q$ se e solo se
	\[
		\sum^{p}_{j=0}a_j = 1\ \ \ \sum^{p}_{j=0}(-j)^ka_j + k \sum^{p}_{j=-1}(-j)^{k-1}b_j=1\ \ \ k=1,\ldots,q
	.\] 
\end{prop}
\begin{dimo}
	$q=1$ (se proprio vogliamo possiamo fare noi i conti per gli altri casi)\\
	Consideriamo lo sviluppo di Taylor di $y_{n-j}=y(t_{n-h})$ in $t_n$, ricordando che  $t_{n-j} = t_n - jh = nh-jh = (n-j)h$
	\[
	h\tau_{n+1}(h) = y_{n+1} - \left[ \sum^{p}_{j=0}a_jy_{n-j} + h \sum^{p}_{j=-1}b_jf(t_{n-j},y_{n-j}) \right] \ \ \ n\geq p\]
	\[
		\begin{split}
			y(t_{n-j})&= y(t_n) + \dot y (t_n)(-jh) + O(h^2)\\
			f(t_{n-j}, y(t_{n-j})) &= f(t_n,y_n) + O(h)
		\end{split}
	.\] 
	Sostituendo nello schema LMM
	\[
		\begin{split}
		&y_{n+1} - \sum^{p}_{j=0}a_j(y_n + \dot y(t_n)(-jh)+O(h^2)) - h \sum^{p}_{j=-1}b_j (f_n + O(h))\\
		&= y_{n+1} - \sum^{p}_{j=0} a_jy_n + h \sum^{p}_{j=0}a_jj\dot y(t_n) - O(h^2) \sum^{p}_{k=0}a_j - hf_n \sum_{j=-1}^{p} - O(h^2) \sum^{p}_{j=-1}b_j
		\end{split}
	.\] 
	Dividiamo tutto per  $h$ per trovare $\tau_{n+1}$
	 \[
		 \begin{split}
		& \tau_{n+1}(h) = \frac{y_{n+1}} h- \sum^{p}_{j=0}\frac{a_jy_n} h + \frac {\cancel{h}}{\cancel{h}}\dot y(t_n) \left( \sum^{p}_{j=0}ja_j - \sum_{j=-1}^{p}b_j \right) + O(h)\\
		&\tau_{n+1}(h) = \frac{y_{n+1}-y_n}h + \frac{y_n - \sum^{p}_{j=0}a_jy_n}h + \dot y (t_n) \left( \sum^{p}_{j=0}ja_j- \sum^{p}_{j=-1}b_j  \right)+ O(h)\\
		&\tau_{n+1} = \dot y (t_n) + \frac {y_n}h(1 - \sum^{p}_{j=0}a_j) + \dot y (t_n) (\ldots ) + O(h)
		\end{split}
	.\] 
	Da cui ricaviamo le condizioni della proposizione.\\
	Quindi 
	\[
		t_{n+1} (h) \xrightarrow{h \rightarrow 0} 0 \Leftrightarrow \sum^{p}_{j=0}a_j = 1 \text { e } \sum^{p}_{j=0}ja_j + \sum^{p}_{j=-1}b_j = -1 
	.\] 
	
\end{dimo}
\textbf{Nota}\\
Nella notazione compatta compare $0^0$, per convenzione lo poniamo uguale a $1$.\\
\textbf{Esempio} midpoint\\
$u_{n+1} = u_{n-1} + 2hf_n\ \ \ p = 1 \Rightarrow $  LMM a $2$ passi con \\
$a_0 = 0, \ a_1 =1, \ b_{-1}=0, \ b_{0}=2,\ b_{1} = 0$\\
Notiamo che  $ \sum^{p}_{j=0}a_j=1$ e $ \sum^{p}_{j=0}ja_j - \sum^{p}_{j=-1}b_j = 0\cdot a_0 + 1\cdot a_1 - b_{-1}- b_0- b_1 = 1 - 0 -2 -0  = -1$\\
Proviamo $k=2$ per vedere se il metodo è di ordine  $2$\\
sostituendo nella formula ottengo $1 + 2\cdot 0 = 1$ quindi è di ordine 2\\
 $k = 3$  \\
 $-1 + 3\cdot 0 = -1$ che non soddisfa la condizione.\\
  \textbf{Esercizio}\\
  Verificare che il metodo di Simpson è di ordine?

\end{document}
